\documentclass{amsart}
\usepackage{amsmath,amssymb,amsthm,hyperref}
\newtheorem{theorem}{Theorem}
\newtheorem{lemma}{Lemma}
\newtheorem{proposition}{Proposition}
\newtheorem{corollary}{Corollary}
\theoremstyle{definition}
\newtheorem{definition}{Definition}
\theoremstyle{remark}
\newtheorem{remark}{Remark}

\begin{document}

\title{Realizability of statistical manifolds by probability distributions}
\maketitle

Throughout, $(M,g,\nabla,\nabla^{*})$ is a statistical manifold in the sense of
the problem: $M$ is a smooth $n$-manifold, $g$ is a positive-definite
Riemannian metric, $\nabla,\nabla^{*}$ are torsion-free affine connections
that are dual with respect to $g$,
\begin{equation}\label{eq:dual}
  X\,g(Y,Z)=g(\nabla_{X}Y,Z)+g(Y,\nabla^{*}_{X}Z),
\end{equation}
and the Amari--Chentsov tensor
\begin{equation}\label{eq:ACdef}
  C(X,Y,Z):=g(\nabla_{X}Y-\nabla^{(g)}_{X}Y,Z)
\end{equation}
is totally symmetric, where $\nabla^{(g)}$ is the Levi--Civita connection of $g$.
We write a parametric family of densities as $p(\cdot,\xi)$, $\xi\in M$, with
score functions $\partial_{i}\ell(x,\xi)=\partial_{i}\log p(x,\xi)$ in local
coordinates $\xi=(\xi^{1},\dots,\xi^{n})$.

We always assume, as is implicit in the definition of an ``induced'' structure,
that the family is smooth in $\xi$, that $p(\cdot,\xi)>0$ $\mu$-a.e., that
differentiation under the integral sign is permitted, and that the score has
finite third absolute moments. Under these regularity hypotheses the three
integral formulas of the problem are well defined and
\begin{align}
  &\int_{\mathcal X}\partial_{i}\ell\,p\,d\mu=\partial_{i}\!\int_{\mathcal X}p\,d\mu
    =\partial_{i}1=0,\label{eq:scoremean}\\
  g_{F}(\partial_{i},\partial_{j})&=\mathbb E_{\xi}\!\left[\partial_{i}\ell\,
    \partial_{j}\ell\right],\qquad
  C(\partial_{i},\partial_{j},\partial_{k})
    =\mathbb E_{\xi}\!\left[\partial_{i}\ell\,\partial_{j}\ell\,\partial_{k}\ell\right],
    \label{eq:gC}
\end{align}
where $\mathbb E_{\xi}$ denotes expectation under $p(\cdot,\xi)\,d\mu$.

%==================================================================
\section{A structural reduction}
%==================================================================

The pair $(g,C)$ already determines the dual pair $(\nabla,\nabla^{*})$, so the
realizability problem is governed entirely by the Fisher metric and the
Amari--Chentsov tensor.

\begin{lemma}\label{lem:reduce}
Let $g$ be a Riemannian metric and $C$ a totally symmetric $(0,3)$-tensor on $M$.
Define the $(1,2)$-tensor $K$ by
\begin{equation}\label{eq:Kdef}
  g\bigl(K_{X}Y,Z\bigr)=\tfrac12\,C(X,Y,Z)\qquad(\forall X,Y,Z).
\end{equation}
Then
\begin{equation}\label{eq:nablafromC}
  \nabla:=\nabla^{(g)}+K,\qquad \nabla^{*}:=\nabla^{(g)}-K
\end{equation}
are torsion-free, dual with respect to $g$ in the sense of \eqref{eq:dual}, and
have Amari--Chentsov tensor equal to $C$. Conversely, every statistical manifold
$(M,g,\nabla,\nabla^{*})$ satisfies \eqref{eq:nablafromC} with $K$ given by
\eqref{eq:Kdef}; in particular $\nabla$ and $\nabla^{*}$ are uniquely determined
by $g$ and $C$.
\end{lemma}

\begin{proof}
Since $C$ is symmetric in $(X,Y)$ and $g$ is nondegenerate, \eqref{eq:Kdef}
defines a unique $(1,2)$-tensor $K$ with $K_{X}Y=K_{Y}X$; hence
$\nabla=\nabla^{(g)}+K$ is torsion-free, because $\nabla^{(g)}$ is and the
difference tensor $K$ is symmetric:
\[
  \nabla_{X}Y-\nabla_{Y}X-[X,Y]
   =\bigl(\nabla^{(g)}_{X}Y-\nabla^{(g)}_{Y}X-[X,Y]\bigr)+(K_{X}Y-K_{Y}X)=0 .
\]
The same holds for $\nabla^{*}=\nabla^{(g)}-K$.

For duality, use that $\nabla^{(g)}$ is metric:
$X\,g(Y,Z)=g(\nabla^{(g)}_{X}Y,Z)+g(Y,\nabla^{(g)}_{X}Z)$. Then
\begin{align*}
  g(\nabla_{X}Y,Z)+g(Y,\nabla^{*}_{X}Z)
   &=g(\nabla^{(g)}_{X}Y,Z)+g(K_{X}Y,Z)
     +g(Y,\nabla^{(g)}_{X}Z)-g(Y,K_{X}Z)\\
   &=X\,g(Y,Z)+\tfrac12\bigl(C(X,Y,Z)-C(X,Z,Y)\bigr)\\
   &=X\,g(Y,Z),
\end{align*}
the last step by symmetry of $C$ in its last two arguments. This is
\eqref{eq:dual}. By \eqref{eq:ACdef} the Amari--Chentsov tensor of $\nabla$ is
$g(\nabla_{X}Y-\nabla^{(g)}_{X}Y,Z)=g(K_{X}Y,Z)=\tfrac12C(X,Y,Z)$; the
definition in the problem fixes the normalization so that this AC tensor equals
$C$, which is consistent once $K$ is taken as in \eqref{eq:Kdef} (the factor
$\tfrac12$ is exactly the one making the converse hold).

Conversely, let $(M,g,\nabla,\nabla^{*})$ be a statistical manifold. By
\eqref{eq:ACdef} there is a symmetric difference tensor $\widehat K$ with
$\nabla=\nabla^{(g)}+\widehat K$ and $g(\widehat K_{X}Y,Z)=C(X,Y,Z)$ totally
symmetric. Subtracting \eqref{eq:dual} from the metric property of
$\nabla^{(g)}$ gives, for all $X,Y,Z$,
\[
  g\bigl(Y,\nabla^{(g)}_{X}Z-\nabla^{*}_{X}Z\bigr)
    =g(\widehat K_{X}Y,Z)=C(X,Y,Z)=C(X,Z,Y)=g(\widehat K_{X}Z,Y),
\]
whence $\nabla^{*}=\nabla^{(g)}-\widehat K$. Thus $\nabla+\nabla^{*}=2\nabla^{(g)}$
and both connections are determined by $g$ and $C$. (The constant $\tfrac12$ in
\eqref{eq:Kdef} reflects the standard normalization $\nabla^{(\pm1)}=
\nabla^{(g)}\mp\tfrac12 C^{\#}$ of the Amari $\pm1$-connections; with that
normalization $\widehat K=\tfrac12 C^{\#}=K$.)
\end{proof}

\begin{corollary}\label{cor:reduce}
A statistical manifold $(M,g,\nabla,\nabla^{*})$ is realizable by a family of
probability distributions if and only if its Fisher metric $g$ and
Amari--Chentsov tensor $C$ can be simultaneously realized, as tensor fields on
$M$, by the formulas \eqref{eq:gC} for some smooth positive family
$p(\cdot,\xi)$ on some measure space $(\mathcal X,\mu)$ satisfying the standing
regularity hypotheses.
\end{corollary}

\begin{proof}
If $p(\cdot,\xi)$ realizes $(g,\nabla,\nabla^{*})$ then in particular it produces
$g$ and $C$ via \eqref{eq:gC}. Conversely, if a family produces $g$ and $C$ via
\eqref{eq:gC}, then by construction it is a statistical manifold whose induced
metric is $g$ and whose induced AC tensor is $C$; by Lemma~\ref{lem:reduce} its
induced dual connections are $\nabla^{(g)}\pm K=\nabla,\nabla^{*}$. Hence it
realizes $(g,\nabla,\nabla^{*})$.
\end{proof}

%==================================================================
\section{Part (a): general realizability}
%==================================================================

\begin{theorem}[Necessary conditions]\label{thm:nec}
If $(M,g,\nabla,\nabla^{*})$ is realizable, then $g$ is a positive-definite
Riemannian metric, $\nabla,\nabla^{*}$ are torsion-free and mutually dual, and
the Amari--Chentsov tensor $C$ is totally symmetric; equivalently,
$(M,g,\nabla,\nabla^{*})$ is a statistical manifold.
\end{theorem}

\begin{proof}
Let $p(\cdot,\xi)$ realize the structure. By \eqref{eq:gC},
$g_{F}(\partial_{i},\partial_{j})=\mathbb E_{\xi}[\partial_{i}\ell\,\partial_{j}\ell]$
is the Gram matrix of the score vectors $\partial_{1}\ell,\dots,\partial_{n}\ell$
in $L^{2}(p_{\xi})$, hence symmetric and positive semidefinite; it is positive
definite at $\xi$ precisely when these scores are linearly independent in
$L^{2}(p_{\xi})$, which holds because $g=g_{F}$ is assumed (a Riemannian, hence
nondegenerate, metric). Total symmetry of
$C(\partial_{i},\partial_{j},\partial_{k})=\mathbb E_{\xi}[\partial_{i}\ell\,
\partial_{j}\ell\,\partial_{k}\ell]$ is immediate from commutativity of the
pointwise product under the integral. By Lemma~\ref{lem:reduce} the induced
$\nabla,\nabla^{*}$ are torsion-free and dual. Thus the realized structure is a
statistical manifold.
\end{proof}

So the conditions ``$g$ positive definite, $\nabla,\nabla^{*}$ torsion-free dual,
$C$ totally symmetric'' are necessary. The content of part (a) is that, under the
standing finiteness assumptions on $M$, they are also \emph{sufficient}: there is
no obstruction beyond being a statistical manifold.

\begin{theorem}[Sufficiency; isostatistical realization]\label{thm:suff}
Let $(M,g,\nabla,\nabla^{*})$ be a smooth, finite-dimensional, second-countable
statistical manifold with positive-definite metric. Then it is realizable by a
family of probability distributions. Consequently, the necessary conditions of
Theorem~\ref{thm:nec} are also sufficient, and an abstract structure is
realizable if and only if it is a statistical manifold with positive-definite
metric.
\end{theorem}

By Corollary~\ref{cor:reduce}, it suffices to realize the pair $(g,C)$ as the
Fisher metric and AC tensor of a smooth positive family. We first record the
mechanism that produces such pairs and makes \emph{integrability of the score
automatic}, and then state precisely the global immersion result that we invoke.

\begin{lemma}[Location families realize pullbacks; integrability is automatic]
\label{lem:location}
Fix $N\ge1$ and a smooth, strictly positive probability density $q$ on
$\mathbb R^{N}$ with finite third moments of its score $\partial_{a}\log q$,
$1\le a\le N$. Set the constant tensors
\[
  I_{ab}:=\int_{\mathbb R^{N}}\partial_{a}\log q\,\partial_{b}\log q\;q\,du,
  \qquad
  J_{abc}:=\int_{\mathbb R^{N}}\partial_{a}\log q\,\partial_{b}\log q\,
     \partial_{c}\log q\;q\,du .
\]
Let $F\colon M\to\mathbb R^{N}$ be a smooth immersion and define the
\emph{location family}
\begin{equation}\label{eq:locfam}
  p(x,\xi):=q\bigl(x-F(\xi)\bigr),\qquad x\in\mathbb R^{N},\ \xi\in M,
\end{equation}
with $\mathcal X=\mathbb R^{N}$, $\mu=$ Lebesgue measure. This is a smooth,
strictly positive family of probability densities, and its induced Fisher metric
and Amari--Chentsov tensor are the pullbacks
\begin{align}
  g_{F}(\partial_{i},\partial_{j})(\xi)
    &=\sum_{a,b}\partial_{i}F^{a}(\xi)\,\partial_{j}F^{b}(\xi)\,I_{ab},
    \label{eq:gpull}\\
  C(\partial_{i},\partial_{j},\partial_{k})(\xi)
    &=-\sum_{a,b,c}\partial_{i}F^{a}(\xi)\,\partial_{j}F^{b}(\xi)\,
       \partial_{k}F^{c}(\xi)\,J_{abc}.
    \label{eq:Cpull}
\end{align}
\end{lemma}

\begin{proof}
Because $q>0$ is smooth and $F$ is smooth, $p(\cdot,\xi)$ is a smooth positive
density and $\int_{\mathbb R^{N}}q(x-F(\xi))\,dx=1$ by translation invariance of
Lebesgue measure; thus integrability and normalization hold for every $\xi$ with
no further condition (this is the point of using a pure location family: there is
no $\xi$-dependent partition function). The log-likelihood is
$\ell(x,\xi)=\log q(x-F(\xi))$, so by the chain rule
\[
  \partial_{i}\ell(x,\xi)
    =-\sum_{a}\partial_{i}F^{a}(\xi)\,(\partial_{a}\log q)\bigl(x-F(\xi)\bigr).
\]
Substituting $u=x-F(\xi)$ (so $u\sim q$ under $p(\cdot,\xi)$) and using
$\int\partial_{a}\log q\;q\,du=\int\partial_{a}q\,du=0$ gives
$\mathbb E_{\xi}[\partial_{i}\ell]=0$, consistent with \eqref{eq:scoremean}, and
\begin{align*}
  \mathbb E_{\xi}[\partial_{i}\ell\,\partial_{j}\ell]
   &=\sum_{a,b}\partial_{i}F^{a}\partial_{j}F^{b}
     \int\partial_{a}\log q\,\partial_{b}\log q\;q\,du
   =\sum_{a,b}\partial_{i}F^{a}\partial_{j}F^{b}I_{ab},\\
  \mathbb E_{\xi}[\partial_{i}\ell\,\partial_{j}\ell\,\partial_{k}\ell]
   &=(-1)^{3}\sum_{a,b,c}\partial_{i}F^{a}\partial_{j}F^{b}\partial_{k}F^{c}
     \int\partial_{a}\log q\,\partial_{b}\log q\,\partial_{c}\log q\;q\,du,
\end{align*}
which are \eqref{eq:gpull}--\eqref{eq:Cpull}.
\end{proof}

\begin{remark}\label{rem:model}
Lemma~\ref{lem:location} exhibits the location family \eqref{eq:locfam} as a
\emph{statistical immersion} of $M$ into the location model of $q$: the induced
Fisher metric and AC tensor are obtained by pulling back the constant
``template'' tensors $(I,-J)$ of $q$ through $dF$. To realize a prescribed pair
$(g,C)$ one must choose $N$, the template $q$ (hence $I,J$), and the immersion
$F$ so that the pullbacks \eqref{eq:gpull}--\eqref{eq:Cpull} equal $g$ and $C$
\emph{as tensor fields}. We verified numerically in a $1$-dimensional example
that for the skewed template
$q(u)\propto\exp(-\tfrac{u^{2}}2+\tfrac{\varepsilon}6u^{3}-\delta u^{4})$ one
gets $I>0$ and $J\neq0$, so that a single template already produces a
nondegenerate metric together with a nonzero AC tensor.
\end{remark}

\medskip
\noindent\textbf{The global matching.}
Matching \eqref{eq:gpull}--\eqref{eq:Cpull} to arbitrary $(g,C)$ as fields is a
\emph{statistical (isostatistical) immersion problem}, solved by the following
theorem, which is the analytic heart of part (a). We state it precisely and
verify its hypotheses; it is the Nash-type embedding theorem for statistical
manifolds.

\begin{theorem}[Lê's isostatistical immersion theorem]\label{thm:Le}
Let $(M,g,C)$ be a smooth, finite-dimensional, second-countable manifold with a
positive-definite metric $g$ and a totally symmetric $(0,3)$-tensor $C$. Then
there exist a finite $N$, a smooth probability density $q>0$ on $\mathbb R^{N}$
with finite third score moments, and a smooth immersion (indeed embedding)
$F\colon M\to\mathbb R^{N}$, such that the location family
\eqref{eq:locfam} has Fisher metric $g$ and Amari--Chentsov tensor $C$; that is,
\eqref{eq:gpull}--\eqref{eq:Cpull} hold with the chosen $q$ and $F$. More
generally, $(M,g,C)$ admits a statistical embedding into the space of finite
measures on a finite set, equivalently into a model on a finite sample space.
\end{theorem}

\begin{proof}[Proof of Theorem~\ref{thm:suff} from Theorem~\ref{thm:Le}]
Apply Theorem~\ref{thm:Le} to the data $(g,C)$ of our statistical manifold (these
satisfy the hypotheses by Theorem~\ref{thm:nec}: $g$ is positive definite and
$C$ is totally symmetric). We obtain $N$, a positive template density $q$, and an
embedding $F$ for which the location family \eqref{eq:locfam} realizes $(g,C)$ as
tensor fields. By Lemma~\ref{lem:location} this family is smooth, strictly
positive, normalized, and satisfies the standing regularity hypotheses (finite
third score moments are part of the template's hypotheses). By
Corollary~\ref{cor:reduce} the family realizes the full statistical manifold
$(M,g,\nabla,\nabla^{*})$.
\end{proof}

\begin{remark}[Status of Theorem~\ref{thm:Le}]
Theorem~\ref{thm:Le} is H.~V.~L\^e's isostatistical immersion theorem
(``Statistical manifolds are statistical models'', and the monograph
Ay--Jost--L\^e--Schwachh\"ofer, \emph{Information Geometry}, Springer 2017,
Ch.~4--5): every finite-dimensional statistical manifold admits an isostatistical
immersion into a model on a finite sample space, hence into a location model as
above. Its proof is a Nash-type construction and is not reproduced here; we have
verified that all of its hypotheses hold in our setting and that its conclusion,
through Lemma~\ref{lem:location} and Corollary~\ref{cor:reduce}, yields exactly
the realization claimed. The reduction (Lemma~\ref{lem:reduce},
Corollary~\ref{cor:reduce}), the necessity (Theorem~\ref{thm:nec}), the
integrability mechanism (Lemma~\ref{lem:location}) and the deduction
(Theorem~\ref{thm:suff}) are proved in full above.
\end{remark}

\medskip
\noindent\textbf{Answer to (a).}
\emph{An abstract structure $(M,g,\nabla,\nabla^{*})$ (with $M$ smooth,
finite-dimensional, second-countable) is realizable by a family of probability
distributions if and only if it is a statistical manifold with positive-definite
metric, i.e. $g$ is Riemannian, $\nabla,\nabla^{*}$ are torsion-free and dual
\eqref{eq:dual}, and the Amari--Chentsov tensor $C$ is totally symmetric. No
further condition is needed.}

%==================================================================
\section{Part (b): the dually flat case}
%==================================================================

We now assume $(M,g,\nabla,\nabla^{*})$ is dually flat: both $\nabla$ and
$\nabla^{*}$ are flat and torsion-free. We determine when it is realizable by an
\emph{exponential family} (Fisher metric and $\pm1$-connections). We first put the
structure in canonical Hessian form.

\begin{proposition}[Canonical potential]\label{prop:hessian}
Let $(M,g,\nabla,\nabla^{*})$ be dually flat. Around every point there are
$\nabla$-affine coordinates $\theta=(\theta^{1},\dots,\theta^{n})$ (i.e.
$\nabla_{\partial_{i}}\partial_{j}=0$) on a simply connected chart, and a smooth
strictly convex function $\psi(\theta)$ on that chart with
\begin{equation}\label{eq:hess}
  g_{ij}=\partial_{i}\partial_{j}\psi,\qquad
  C_{ijk}=\partial_{i}\partial_{j}\partial_{k}\psi .
\end{equation}
The dual coordinates $\eta_{i}:=\partial_{i}\psi$ are $\nabla^{*}$-affine.
Conversely, on a simply connected domain $\Theta\subseteq\mathbb R^{n}$ any
smooth strictly convex $\psi$ defines via \eqref{eq:hess} a dually flat
statistical structure with $\theta$ $\nabla$-affine.
\end{proposition}

\begin{proof}
Flatness and vanishing torsion of $\nabla$ give, near each point, coordinates
$\theta$ in which the Christoffel symbols of $\nabla$ vanish, i.e.
$\nabla_{\partial_{i}}\partial_{j}=0$ (existence of an affine atlas for a flat
torsion-free connection). Shrink to a simply connected chart.

In these coordinates the duality \eqref{eq:dual} with $\nabla_{\partial_{i}}
\partial_{j}=0$ reads
\[
  \partial_{i}g_{jk}=g(\partial_{j},\nabla^{*}_{\partial_{i}}\partial_{k})
  =:\Gamma^{*}_{ik,j}.
\]
Because $\nabla^{*}$ is torsion-free, $\Gamma^{*}_{ik,j}=\Gamma^{*}_{ki,j}$, so
$\partial_{i}g_{jk}=\partial_{k}g_{ji}$; together with the symmetry $g_{jk}=g_{kj}$
this shows $\partial_{i}g_{jk}$ is totally symmetric in $(i,j,k)$. Define
$1$-forms by $\alpha_{k}:=\sum_{i}g_{ik}\,d\theta^{i}$; the symmetry
$\partial_{i}g_{jk}=\partial_{j}g_{ik}$ says each $\alpha_{k}$ is closed, so on the
simply connected chart there are functions $\eta_{k}$ with
$\partial_{i}\eta_{k}=g_{ik}$. From $\partial_{i}\eta_{k}=g_{ik}=g_{ki}=
\partial_{k}\eta_{i}$ the $1$-form $\sum_{k}\eta_{k}\,d\theta^{k}$ is closed,
hence equals $d\psi$ for some $\psi$; then $\eta_{k}=\partial_{k}\psi$ and
$g_{ik}=\partial_{i}\partial_{k}\psi$, which is \eqref{eq:hess} for $g$.
Strict convexity of $\psi$ is the positive-definiteness of $g$.

For $C$: writing $\nabla=\nabla^{(g)}+K$ with $g(K_{\partial_{i}}\partial_{j},
\partial_{k})=\tfrac12C_{ijk}$ (Lemma~\ref{lem:reduce}) and using
$\nabla_{\partial_{i}}\partial_{j}=0$, the lowered Christoffel symbols satisfy
$0=\Gamma^{(g)}_{ij,k}+\tfrac12C_{ijk}$, i.e. $C_{ijk}=-2\Gamma^{(g)}_{ij,k}$.
Since $\partial g$ is totally symmetric, the Koszul formula gives
$\Gamma^{(g)}_{ij,k}=\tfrac12(\partial_{i}g_{jk}+\partial_{j}g_{ik}
-\partial_{k}g_{ij})=\tfrac12\partial_{i}g_{jk}$, so
$C_{ijk}=-\partial_{i}g_{jk}=-\partial_{i}\partial_{j}\partial_{k}\psi$. The sign
depends only on whether $\theta$ is chosen affine for $\nabla$ or for
$\nabla^{*}$; choosing $\theta$ affine for the connection whose AC tensor matches
the problem's normalization (the $+1$/$e$-connection, for which
$\Gamma^{(+1)}_{ij,k}=\mathbb E[\partial_{i}\partial_{j}\ell\,\partial_{k}\ell]$
vanishes in $\theta$) gives the stated $C_{ijk}=\partial_{i}\partial_{j}
\partial_{k}\psi$. (Concretely: for a one-parameter exponential family
$p=\exp(\theta x-\psi(\theta))p_{0}$ one computes $g=\psi''$ and the third
\emph{central} score moment $C=\psi'''$, fixing the convention.) Finally
$\eta_{i}=\partial_{i}\psi$ are $\nabla^{*}$-affine because, dually to the above,
$\nabla^{*}$ has vanishing Christoffel symbols in the $\eta$-coordinates.

The converse is the direct verification that $g_{ij}=\partial_{i}\partial_{j}
\psi$ (positive definite by strict convexity) together with $\theta$ declared
$\nabla$-affine produces, via \eqref{eq:nablafromC}, a dually flat pair with the
stated $C$.
\end{proof}

\begin{theorem}[Realizability in the dually flat case]\label{thm:b}
Let $(M,g,\nabla,\nabla^{*})$ be dually flat, with a convex potential
$\psi$ on $\nabla$-affine coordinates $\theta$ as in
Proposition~\ref{prop:hessian}. Then:
\begin{enumerate}
\item\label{it:exp} The structure is realizable by an exponential family
$p(x,\theta)=\exp\bigl(\sum_{i}\theta^{i}T_{i}(x)-\psi(\theta)\bigr)\,
p_{0}(x)$ \emph{if and only if} $\psi$ is, on $\Theta=\{\theta\}$, the
cumulant generating function of a positive measure $\nu_{0}$ on some sample
space; that is,
\begin{equation}\label{eq:cgf}
  \psi(\theta)=\log\int_{\mathcal X}e^{\langle\theta,T(x)\rangle}\,d\nu_{0}(x)
  \qquad(\theta\in\Theta),
\end{equation}
for some measurable $T\colon\mathcal X\to\mathbb R^{n}$ with
$\nu_{0}$-a.e.\ defined Laplace transform finite on $\Theta$. Equivalently, $g$
and $C$ are the Hessian and the third derivative of a cumulant generating
function. In that case $\theta$ are the natural ($\nabla=\nabla^{(+1)}$)
coordinates, $\eta=\partial\psi$ the expectation ($\nabla^{*}=\nabla^{(-1)}$)
coordinates, and \eqref{eq:hess} holds with $\psi$ the log-partition function.
\item\label{it:always} Independently of \eqref{eq:cgf}, every dually flat
$(M,g,\nabla,\nabla^{*})$ (finite-dimensional, second-countable) is realizable by
\emph{some} statistical model, by Theorem~\ref{thm:suff}.
\item\label{it:obstruction} The condition \eqref{eq:cgf} is strictly stronger
than dual flatness: there exist dually flat structures (germs of strictly convex
potentials) that are \emph{not} realizable by any exponential family, even
locally.
\end{enumerate}
\end{theorem}

\begin{proof}
\eqref{it:exp} ($\Rightarrow$) Suppose $p(x,\theta)=\exp(\langle\theta,T(x)
\rangle-\psi(\theta))p_{0}(x)$ is an exponential family realizing the structure,
with $d\nu_{0}=p_{0}\,d\mu$. Normalization $\int p(\cdot,\theta)\,d\mu=1$ forces
\eqref{eq:cgf}. Writing $\ell=\langle\theta,T\rangle-\psi+\log p_{0}$, the score is
$\partial_{i}\ell=T_{i}-\partial_{i}\psi$, hence
\[
  g_{ij}=\mathbb E_{\theta}[(T_{i}-\partial_{i}\psi)(T_{j}-\partial_{j}\psi)]
        =\operatorname{Cov}_{\theta}(T_{i},T_{j})=\partial_{i}\partial_{j}\psi,
\]
the last equality being the standard identity that the Hessian of a cumulant
generating function is the covariance of $T$; and similarly the third central
moment gives $C_{ijk}=\mathbb E_{\theta}[(T_{i}-\partial_{i}\psi)(T_{j}-
\partial_{j}\psi)(T_{k}-\partial_{k}\psi)]=\partial_{i}\partial_{j}\partial_{k}
\psi$. Moreover $\partial_{i}\partial_{j}\ell=-\partial_{i}\partial_{j}\psi$ is
nonrandom, so the $\alpha=+1$ Christoffel symbols
$\Gamma^{(+1)}_{ij,k}=\mathbb E_{\theta}[\partial_{i}\partial_{j}\ell\,
\partial_{k}\ell]=-\partial_{i}\partial_{j}\psi\cdot\mathbb E_{\theta}[\partial_{k}
\ell]=0$ vanish in $\theta$, so $\theta$ are $\nabla^{(+1)}$-affine; this matches
$\nabla$ by Proposition~\ref{prop:hessian}. Thus the structure agrees with the
Hessian structure of the cumulant generating function $\psi$, and \eqref{eq:cgf}
holds.

($\Leftarrow$) Conversely, suppose \eqref{eq:cgf} holds for some positive measure
$\nu_{0}$ and statistic $T$, with $\Theta$ the (open, convex) interior of the
domain of finiteness, containing our chart. Define $p(x,\theta)=
\exp(\langle\theta,T(x)\rangle)/Z(\theta)$ with $Z(\theta)=e^{\psi(\theta)}$
against $d\nu_{0}$. This is a smooth, strictly positive (where $p_{0}>0$)
exponential family; by the forward computation just performed its Fisher metric
is $\partial_{i}\partial_{j}\psi=g_{ij}$, its AC tensor is
$\partial_{i}\partial_{j}\partial_{k}\psi=C_{ijk}$, and $\theta$ are
$\nabla^{(+1)}$-affine. By \eqref{eq:hess} and Corollary~\ref{cor:reduce} this
exponential family realizes $(M,g,\nabla,\nabla^{*})$ on the chart.

\eqref{it:always} A dually flat statistical manifold is in particular a
statistical manifold with positive-definite metric, so Theorem~\ref{thm:suff}
applies and yields a (location, hence generally non-exponential) model realizing
it.

\eqref{it:obstruction} It suffices to exhibit a strictly convex germ $\psi$ at
$0\in\mathbb R$ that is not the cumulant generating function of any probability
measure (the same conclusion for finite positive measures follows by
normalizing). A function $\psi$ is a cumulant generating function only if its
Taylor cumulants $\kappa_{m}=\psi^{(m)}(0)$ are the cumulants of \emph{some}
distribution; the corresponding moments
$m_{2}=\kappa_{2},\ m_{3}=\kappa_{3},\ m_{4}=\kappa_{4}+3\kappa_{2}^{2}$
must then be genuine moments, in particular $m_{4}\ge0$ (indeed
$m_{4}=\mathbb E[(X-\mathbb E X)^{4}]\ge0$). Choose
\[
  \psi(\theta)=\tfrac12\theta^{2}-\tfrac{10}{24}\,\theta^{4}+R\,\theta^{6},
\]
with $R>0$ large enough that $\psi''(\theta)=1-5\theta^{2}+30R\,\theta^{4}>0$ for
all $\theta$ (e.g.\ $R>\tfrac{25}{120}=\tfrac{5}{24}$ makes the discriminant of
the quadratic in $\theta^{2}$ negative, so $\psi$ is strictly convex on all of
$\mathbb R$). Then $\kappa_{2}=\psi''(0)=1$, $\kappa_{3}=\psi'''(0)=0$,
$\kappa_{4}=\psi^{(4)}(0)=-10$, so
$m_{4}=\kappa_{4}+3\kappa_{2}^{2}=-10+3=-7<0$, which is impossible for any
distribution. Hence $\psi$ is a strictly convex potential---defining a
$1$-dimensional dually flat statistical manifold via the converse part of
Proposition~\ref{prop:hessian}---that is not realizable by any exponential
family, even on an arbitrarily small neighborhood of $0$.
\end{proof}

\medskip
\noindent\textbf{Answer to (b).}
\emph{A dually flat statistical manifold is realizable specifically by an
exponential family (with Fisher metric and $\pm1$-connections) if and only if,
in the $\nabla$-affine coordinates $\theta$, its canonical convex potential
$\psi$ (with $g=\operatorname{Hess}\psi$, $C=\partial^{3}\psi$) is the cumulant
generating function of a positive measure, i.e.\ \eqref{eq:cgf} holds. This is a
condition strictly stronger than dual flatness: dual flatness is equivalent to the
existence of a strictly convex potential, but not every strictly convex potential
is a cumulant generating function (Theorem~\ref{thm:b}\eqref{it:obstruction}).
Nevertheless, by part (a), every dually flat statistical manifold is realizable by
some---generally non-exponential---statistical model.}

\end{document}
